\subsection{Localized acceptance}\label{sec:localized}
The acceptance rule of Proposition~\ref{prop:accept} needs a certified gap
below $1/(\Omega W)$, as small as about $2^{-315}$ on the instances of
experiment E4 (Section~\ref{sec:comp-exact}), which have at most six
variables. The filtering
history permits a different test that accepts a candidate as soon as the retained box
is small compared with a first-order margin.

Consider stage $j$ of a filtering run, such as a trial of CT: a grid $G$ for
the current box $X^{(j)}$ with corrections $d_i(v)=L_iw_i(v)^2/8$ and the
grids $G_i=\{\ell_i,u_i\}$ for $i\notin P$, min-marginals $m_i$, a corrected
minimizer $y_j$, the incumbent value $U$, and the filtered box $X^{(j+1)}$.
The boxes are reached from $X$ by filtering steps whose thresholds are values
of feasible points and at least $U$, as in a trial, where the incumbent value
never increases. So by Proposition~\ref{prop:filter} every $x\in X$ with
$F(x)\le U$ lies in $X^{(j+1)}$, and $y_j\in X^{(j+1)}$. Filtering never
narrows a coordinate $i\notin P$, so $X^{(j+1)}_i=X^{(j)}_i=X_i$.

\begin{lemma}[Node exclusion]\label{lem:node}
Let $x\in X^{(j)}$ with $F(x)\le U$.
\begin{enumerate}[label=(\alph*)]
\item If $i\in I_Z$, then $x_i$ lies in the set $Z_i$ of nodes $v$ with
$m_i(v)\le U$ and of integers strictly inside grid intervals $[a,a']$ with
$a'-a\ge2$ and $\min\{m_i(a),m_i(a')\}\le U$.
\item If $i\notin P$, $G_i=\{a,a'\}$ and $m_i(a)>U$, then setting $x_i=a'$
gives a point of $X^{(j)}$ with value at most $F(x)$.
\end{enumerate}
\end{lemma}

\begin{proof}
(a) If $x_i$ is a node $v$, Proposition~\ref{prop:cellwise}(c) gives
$m_i(v)\le m_i(v)+d_i(v)\le F(x)\le U$. Otherwise $x_i$ is an integer strictly
inside a grid interval $[a,a']$, which then has $a'-a\ge2$, and
Proposition~\ref{prop:cellwise}(c) gives $\min\{m_i(a),m_i(a')\}\le F(x)\le U$.
(b) Since $L_i=0$, $F$ is concave along coordinate $i$, so $F(x)$ is at least
the smaller of its values at $x_i=a$ and $x_i=a'$. By
Proposition~\ref{prop:cellwise}(c) the first is at least $m_i(a)>U\ge F(x)$,
so the second is at most $F(x)$.
\end{proof}

\begin{proposition}[Localized exact acceptance]\label{prop:local}
Define the narrowed box $\tilde X\subseteq X^{(j+1)}$ coordinatewise: for
$i\in I_Z\cap P$ let $\tilde X_i$ be the integer hull of
$Z_i\cap X^{(j+1)}_i$; for $i\notin P$ and $G_i=\{a,a'\}$ let
$\tilde X_i=\{a'\}$ if $m_i(a)>U\ge m_i(a')$, $\tilde X_i=\{a\}$ if
$m_i(a')>U\ge m_i(a)$, and $\tilde X_i=X_i$ otherwise; and let
$\tilde X_i=X^{(j+1)}_i$ for $i\in I_C\cap P$. Let $\tilde x\in\tilde X$
with $F(\tilde x)\le U$, $\zeta=\nabla F(\tilde x)$,
$J_+=\{i:\tilde X_i\text{ has positive length}\}$, and
$\tilde X_i=[\alpha_i,\alpha_i']$ (or its integer points) for $i\in J_+$.
Suppose every $i\in J_+$ falls under one of the cases
\[
\begin{array}{lll}
\text{(i)}&\tilde x_i=\alpha_i,\ \zeta_i\ge0:&\mu_i=\zeta_i/(\alpha_i'-\alpha_i);\\
\text{(ii)}&\tilde x_i=\alpha_i',\ \zeta_i\le0:&\mu_i=-\zeta_i/(\alpha_i'-\alpha_i);\\
\text{(iii)}&\alpha_i<\tilde x_i<\alpha_i',\ \zeta_i=0:&\mu_i=0;\\
\text{(iv)}&i\in I_Z\text{ otherwise}:&\mu_i=-\abs{\zeta_i},
\end{array}
\]
and that $H_{J_+J_+}+2\diag(\mu_{J_+})\succeq0$. Then $\tilde x$ is a global
minimizer.
\end{proposition}

\begin{proof}
Let $x\in X$. If $F(x)>U$, then $F(x)>F(\tilde x)$. Otherwise
$x\in X^{(j+1)}$. Applying Lemma~\ref{lem:node}(b) to the coordinates
$i\notin P$ whose $\tilde X_i$ is a single endpoint keeps the point in
$X^{(j+1)}$ and does not increase $F$, and gives a point $x''$ that, by
Lemma~\ref{lem:node}(a), lies in $\tilde X$. Let $d=x''-\tilde x$; then
$d_i=0$ for $i\notin J_+$ and $\alpha_i-\tilde x_i\le d_i\le\alpha_i'-\tilde x_i$
for $i\in J_+$. In cases (i) and (ii),
$\zeta_id_i=\abs{\zeta_i}\abs{d_i}\ge\mu_id_i^2$; in case (iii) both sides
vanish; in case (iv) $d_i\in\Z$, so
$\zeta_id_i\ge-\abs{\zeta_i}\abs{d_i}\ge-\abs{\zeta_i}d_i^2$. Hence
$F(x'')-F(\tilde x)=\zeta^\T d+\frac12d^\T Hd\ge\frac12d_{J_+}^\T\bigl(H_{J_+J_+}
+2\diag(\mu_{J_+})\bigr)d_{J_+}\ge0$, and $F(x)\ge F(x'')\ge F(\tilde x)$.
\end{proof}

The test needs one exact $LDL^\T$ factorization of a
$\abs{J_+}\times\abs{J_+}$ matrix, which a checker repeats after replaying the
filtering history. Its validity needs neither uniqueness nor growth, and the
candidate $\tilde x$ is arbitrary. The corrected minimizer lies in
$\tilde X$: for $i\in I_Z\cap P$, $(y_j)_i$ is a node with
$m_i((y_j)_i)\le Q(y_j)=\beta_j\le\OPT\le U$, and for $i\notin P$ the same
inequality excludes the endpoint whose min-marginal exceeds $U$. The
following rule computes a candidate from $\tilde X$ and $y_j$.

\begin{definition}[Face candidate]\label{def:facecand}
Put $\tilde x_i=(y_j)_i$ for $i\in I_Z$ and for $i\notin J_+$. For
$i\in I_C\cap J_+$, put $\tilde x_i=\ell_i$ if $\ell_i\in\tilde X_i$, else
$\tilde x_i=u_i$ if $u_i\in\tilde X_i$, and otherwise call $i$ free. If the
set $J_f$ of free coordinates has $H_{J_fJ_f}$ nonsingular, the \emph{face
candidate} is the point $\tilde x$ whose free part solves
$\nabla_{J_f}F(\tilde x)=0$; otherwise stage $j$ has no face candidate.
\end{definition}

The face candidate costs one exact solve of a $\abs{J_f}\times\abs{J_f}$ linear
system. It is tested by Proposition~\ref{prop:local}, whose hypotheses
include $\tilde x\in\tilde X$ and $F(\tilde x)\le U$.

\begin{corollary}[Acceptance of the face candidate]\label{cor:local}
Let $F$ have point growth with constant $g$ at $x^*$. Let $J_0=J_0(x^*)$ and
$A=J_\partial\cap I_C^+$, the continuous coordinates with $L_i>0$ at which
$x^*$ is at a bound, and assume strict complementarity,
$\partial_iF(x^*)\ne0$ for $i\in A$. Put
$\lambda_A=\min_{i\in A}\abs{\partial_iF(x^*)}$ and
$\Gamma=\norm{H_{AA}}_2+\norm{H_{J_0A}}_2^2/(2g)$,
and let $\delta_X$ be the smallest distance $\abs{x^*_i-e}$ over coordinates
$i\in I_C\cup(I_Z\setminus P)$ and endpoints $e\in\{\ell_i,u_i\}$ with
$e\ne x^*_i$. Put $h^*=\omega^*/\sqrt{n\kappa}$, where $\omega^*$ is the
smallest of the numbers $\frac12$ (if $I_Z\cap P\ne\emptyset$), $\delta_X/6$
(if $I_C\cup(I_Z\setminus P)\ne\emptyset$) and $\lambda_A/(5\Gamma)$ (if
$A\ne\emptyset$); at least one of the first two conditions holds, because
$I_Z\cap P$ and $I_C\cup(I_Z\setminus P)$ cover $[n]$. Consider a trial
of CT with a common mesh (Lemma~\ref{lem:commonmesh}) with
$8\kappa\theta^2\le1$. At every stage $j$ with $h_j\le h^*$ that the trial
reaches and ends by filtering, the face candidate exists, equals $x^*$, and
passes the test of Proposition~\ref{prop:local}. Since $\delta_X\ge1/R$,
$\lambda_A\ge1/(\Delta R)$ and $\log_2\Gamma\le\poly(I)+\log_2\kappa$, the
least $j$ with $h_j\le h^*$ is at most $\poly(I)+O(\log\kappa)$.
\end{corollary}

The proof is in Appendix~\ref{app:localized}. It also shows that at these
stages $\tilde X_i=\{x^*_i\}$ for every integer coordinate and every
coordinate with $L_i=0$, and that $J_f=J_0$. With several minimizers the test
can fail, for example when two optimal vertices differ in a coordinate
$i\notin P$.

In experiment S1 (Section~\ref{sec:comp-localized}) the test accepted the
face candidate within nine stages on 29 of 30 random mixed-integer instances;
on the five of them on which EX used the most stages, one CT run needed 40 to
72 stages to reach the threshold of Proposition~\ref{prop:accept} with the
constants of~\eqref{eq:exact-constants}. The idea of certifying a neighborhood of a
candidate by local information is that of exclusion boxes in interval
branch-and-bound \cite{Neumaier2004}; here it is carried out exactly on the
box produced by the filtering history.
